% ==============================================================================
% CONFIGURATION: PREAMBLE & DANH MỤC GÓI MỞ RỘNG
% ==============================================================================

% extreport has no native 13pt option; set all relative font sizes explicitly.
\usepackage[fontsize=13pt]{fontsize}
\usepackage{ifxetex,ifluatex}
\newif\ifxetexorluatex
\ifxetex
  \xetexorluatextrue
\else
  \ifluatex
    \xetexorluatextrue
  \else
    \xetexorluatexfalse
  \fi
\fi

% --- Cấu hình Font chữ Times New Roman và Tiếng Việt ---
\ifxetexorluatex
  \usepackage{fontspec}
  \IfFileExists{fonts/times.ttf}{%
    \setmainfont{times.ttf}[
      Path = fonts/,
      BoldFont = timesbd.ttf,
      ItalicFont = timesi.ttf,
      BoldItalicFont = timesbi.ttf
    ]
  }{%
    \IfFileExists{font/times.ttf}{%
      \setmainfont{times.ttf}[
        Path = font/,
        BoldFont = timesbd.ttf,
        ItalicFont = timesi.ttf,
        BoldItalicFont = timesbi.ttf
      ]
    }{%
      \setmainfont{Times New Roman}[
        Path = ,
        Extension = ,
        UprightFont = *,
        BoldFont = * Bold,
        ItalicFont = * Italic,
        BoldItalicFont = * Bold Italic
      ]
    }
  }
  \setsansfont{Arial}
  \setmonofont{Courier New}[Scale=0.92]
\else
  \usepackage[utf8]{inputenc}
  \usepackage[T5,T1]{fontenc}
  \usepackage{mathptmx}
\fi

% --- Cấu hình Trang in & Căn lề chuẩn Luận văn / BTL ---
\usepackage[
  a4paper,
  top=25mm,
  bottom=25mm,
  left=30mm,
  right=20mm,
  headheight=20pt,
  headsep=12pt,
  footskip=12mm
]{geometry}

% --- Giãn dòng và thụt đầu dòng ---
\usepackage{setspace}
\usepackage{indentfirst} % Thụt đầu dòng đoạn đầu tiên sau tiêu đề mục chuẩn UTH
\setstretch{1.3} % Giãn dòng 1.3 chuẩn bài tập lớn UTH
\setlength{\parindent}{1.0cm}
\setlength{\parskip}{4pt plus 2pt minus 1pt}
\hyphenation{abstract introduction chunks}
\widowpenalty=10000
\clubpenalty=10000

% --- Toán học nâng cao & Ký hiệu ---
\usepackage{amsmath,amssymb,amsfonts,mathtools,amsthm}
\usepackage{bm}

% --- Đồ họa & Màu sắc ---
\usepackage{graphicx}
\usepackage[table,xcdraw,dvipsnames]{xcolor}
\definecolor{UTHBlue}{HTML}{0054A6}
\definecolor{UTHTeal}{HTML}{008080}
\definecolor{UTHRed}{HTML}{C8102E}
\definecolor{DarkNavy}{HTML}{0F172A}
\definecolor{LightGray}{HTML}{F8FAFC}
\definecolor{BorderGray}{HTML}{CBD5E1}
\definecolor{CodeBg}{HTML}{F1F5F9}
\colorlet{Teal}{UTHTeal}

% --- Bảng biểu chuyên nghiệp ---
\usepackage{booktabs}
\usepackage{tabularx}
\usepackage{multirow}
\usepackage{makecell}
\usepackage{array}
\usepackage{longtable}
\usepackage{needspace}
\usepackage{placeins}
\newcolumntype{Y}{>{\centering\arraybackslash}X}
\newcolumntype{L}{>{\raggedright\arraybackslash}X}

% --- TikZ & Vẽ sơ đồ Vector ---
\usepackage{tikz}
\usetikzlibrary{calc,shapes,arrows.meta,positioning,shadows,fit,backgrounds,trees}

% --- Code Listing & Hộp chú thích ---
\usepackage{listings}
\lstset{
  backgroundcolor=\color{CodeBg},
  basicstyle=\ttfamily\fontsize{9.5}{11.5}\selectfont,
  columns=fullflexible,
  keepspaces=true,
  breakatwhitespace=false,
  xleftmargin=1em,
  xrightmargin=0.2em,
  framexleftmargin=0.4em,
  breaklines=true,
  frame=single,
  rulecolor=\color{BorderGray},
  keywordstyle=\color{UTHBlue}\bfseries,
  commentstyle=\color{gray}\itshape,
  stringstyle=\color{OliveGreen},
  numbers=left,
  numberstyle=\tiny\color{gray},
  numbersep=8pt,
  tabsize=2,
  showstringspaces=false,
  captionpos=b
}

\usepackage[most]{tcolorbox}
\newtcolorbox{calloutbox}[2][]{
  enhanced,
  colback=LightGray,
  colframe=UTHBlue,
  coltitle=white,
  fonttitle=\bfseries\small,
  title={#2},
  boxrule=0.8pt,
  arc=3mm,
  top=3mm, bottom=3mm, left=4mm, right=4mm,
  #1
}

% --- Định dạng Đề mục (Section / Chapter) Chuẩn hóa ---
\usepackage{titlesec}
\renewcommand{\chaptertitlename}{CHƯƠNG}

% Tiêu đề chương có đánh số: CHƯƠNG X (dòng 1), TÊN CHƯƠNG (dòng 2) - IN HOA, MÀU ĐEN, IN ĐẬM, CĂN GIỮA
\titleformat{\chapter}[display]
  {\normalfont\bfseries\fontsize{16pt}{20pt}\selectfont\color{black}\centering}
  {\MakeUppercase{\chaptertitlename\ \thechapter}}
  {0.6ex}
  {\MakeUppercase}
\titlespacing*{\chapter}{0pt}{-40pt}{18pt}

% Tiêu đề chương không đánh số (Tóm tắt, Abstract, Mục lục...): IN HOA, MÀU ĐEN, IN ĐẬM, CĂN GIỮA
\titleformat{name=\chapter,numberless}[block]
  {\normalfont\bfseries\fontsize{16pt}{20pt}\selectfont\color{black}\centering}
  {}
  {0pt}
  {\MakeUppercase}
\titlespacing*{name=\chapter,numberless}{0pt}{-40pt}{18pt}

% Tiêu đề Section cấp 1: cỡ chữ 14pt in đậm theo chuẩn UTH Rule 4.2 (giảm khoảng cách cách trên)
\titleformat{\section}
  {\normalfont\bfseries\fontsize{14pt}{18pt}\selectfont\color{black}}
  {\thesection}{1em}{}
\titlespacing*{\section}{0pt}{10pt}{4pt}

% Tiêu đề Subsection cấp 2: cỡ chữ 13pt in đậm theo chuẩn UTH Rule 4.2
\titleformat{\subsection}
  {\normalfont\bfseries\fontsize{13pt}{16pt}\selectfont\color{black}}
  {\thesubsection}{1em}{}
\titlespacing*{\subsection}{0pt}{8pt}{3pt}

% Tiêu đề Subsubsection cấp 3: cỡ chữ 13pt in nghiêng theo chuẩn UTH Rule 4.2
\titleformat{\subsubsection}
  {\normalfont\itshape\fontsize{13pt}{16pt}\selectfont\color{black}}
  {\thesubsubsection}{1em}{}
\titlespacing*{\subsubsection}{0pt}{6pt}{2pt}

% --- Header & Footer (Cấu hình phân vùng Trang Chuẩn UTH) ---
\usepackage{fancyhdr}

% 1. Style cho Phần đầu (Front Matter - Từ Mục lục trở lên): Số trang TRÊN và CHÍNH GIỮA (14pt)
\fancypagestyle{frontmatterstyle}{%
  \fancyhf{}%
  \fancyhead[C]{\fontsize{14pt}{17pt}\selectfont\thepage}%
  \fancyhead[L]{}%
  \fancyhead[R]{}%
  \fancyfoot{}%
  \renewcommand{\headrulewidth}{0pt}%
  \renewcommand{\footrulewidth}{0pt}%
}

% 2. Style cho Nội dung chính (Main Matter - Từ dưới Mục lục xuống): Số trang DƯỚI và CHÍNH GIỮA (14pt), Header TRÊN BÊN TRÁI (14pt)
\fancypagestyle{mainmatterstyle}{%
  \fancyhf{}%
  \fancyhead[L]{\fontsize{14pt}{17pt}\selectfont\itshape Khai thác Dữ liệu Khoa học \& Hệ thống RAG (UTH)}%
  \fancyhead[C]{}%
  \fancyhead[R]{}%
  \fancyfoot[C]{\fontsize{14pt}{17pt}\selectfont\thepage}%
  \renewcommand{\headrulewidth}{0.4pt}%
  \renewcommand{\footrulewidth}{0pt}%
}

% 3. Style Plain mặc định
\fancypagestyle{plain}{%
  \fancyhf{}%
  \fancyfoot[C]{\fontsize{14pt}{17pt}\selectfont\thepage}%
  \renewcommand{\headrulewidth}{0pt}%
  \renewcommand{\footrulewidth}{0pt}%
}

% --- Hỗ trợ ngắt dòng đường dẫn và URL tự động ---
\usepackage{xurl}
\setlength{\emergencystretch}{2em}
\Urlmuskip=0mu plus 1mu
\newcommand{\ReportNote}[1]{\par\smallskip{\footnotesize\raggedright #1\par}}
\renewcommand{\arraystretch}{1.15}

% --- Hyperlinks & Metadata PDF ---
\usepackage{hyperref}
\hypersetup{
  colorlinks=true,
  linkcolor=black,
  citecolor=ForestGreen,
  urlcolor=BlueViolet,
  pdfborder={0 0 0},
  pdftitle={BTL Khai thác Dữ liệu - UTH},
  pdfauthor={Nhóm 1 - UTH AI/DS}
}

% --- Chú thích hình & bảng ---
\usepackage{caption}
\usepackage{subcaption}
\captionsetup{
  font={small},
  labelfont={bf},
  labelsep=colon,
  justification=centering
}

% --- Chuẩn hóa nhãn tiếng Việt cho Bảng & Hình vẽ ---
\renewcommand{\tablename}{Bảng}
\renewcommand{\figurename}{Hình}
\renewcommand{\bibname}{TÀI LIỆU THAM KHẢO}
\renewcommand{\lstlistingname}{Đoạn mã}

\usepackage{enumitem}
\setlist{itemsep=2pt,topsep=2pt,parsep=0pt}

% --- Cấu hình Mục lục (TOC), Danh mục hình ảnh (LOF), Danh mục bảng biểu (LOT) ---
\usepackage[titles]{tocloft}
% Cùng cơ chế chapter* cho mọi tiêu đề; số mục nằm gần nội dung.
\cftsetindents{chapter}{0em}{1.2em}
\cftsetindents{section}{1.2em}{2.0em}
\cftsetindents{subsection}{3.2em}{2.8em}
\cftsetindents{subsubsection}{6.0em}{3.6em}
\setlength{\cftfignumwidth}{2.2em}
\setlength{\cfttabnumwidth}{2.2em}

\renewcommand{\contentsname}{MỤC LỤC}
\renewcommand{\cftaftertoctitle}{\hfill}

% Thu nhỏ kích thước dấu chấm nối dòng và khoảng cách trong Mục lục & Danh mục
\renewcommand{\cftdot}{\scalebox{0.6}{.}}
\renewcommand{\cftdotsep}{2.5}
\renewcommand{\cftchapdotsep}{2.5}
\renewcommand{\cftsecdotsep}{2.5}
\renewcommand{\cftsubsecdotsep}{2.5}
\renewcommand{\cftsubsubsecdotsep}{2.5}
\renewcommand{\cftfigdotsep}{2.5}
\renewcommand{\cfttabdotsep}{2.5}

% Đường chấm nối dài tự động cho toàn bộ các cấp đề mục
\renewcommand{\cftchapleader}{\cftdotfill{\cftsecdotsep}}
\renewcommand{\cftsecleader}{\cftdotfill{\cftsecdotsep}}
\renewcommand{\cftsubsecleader}{\cftdotfill{\cftsecdotsep}}
\renewcommand{\cftsubsubsecleader}{\cftdotfill{\cftsecdotsep}}

% Đồng bộ Danh mục hình vẽ
\renewcommand{\listfigurename}{DANH MỤC HÌNH VẼ}
\renewcommand{\cftafterloftitle}{\hfill}
\renewcommand{\cftfigleader}{\cftdotfill{\cftfigdotsep}}

% Đồng bộ Danh mục bảng biểu
\renewcommand{\listtablename}{DANH MỤC BẢNG BIỂU}
\renewcommand{\cftafterlottitle}{\hfill}
\renewcommand{\cfttableader}{\cftdotfill{\cfttabdotsep}}

\usepackage{tocbibind} % Đưa mục tài liệu tham khảo vào TOC


% Không kéo giãn dọc trang hoặc căn giữa các trang chỉ có hình.
\raggedbottom
\makeatletter
\setlength{\@fptop}{0pt}
\setlength{\@fpsep}{12pt}
\setlength{\@fpbot}{0pt plus 1fil}
\makeatother
